\Article{25}
Bevor wir weitergehen, wollen wir die geometrische Bedeutung der gewonnenen Ausdrücke verdeutlichen. Um die verschiedenen Richtungen anschaulicher zu machen, folgen wir dem Verfahren, das wir in Disquiss. gen. circa superficies curvas eingeführt haben. Wir ordnen ihnen Punkte auf einer um einen beliebigen Mittelpunkt beschriebenen Kugeloberfläche vom Radius $=1$ zu. Zuerst bezeichnen wir also die Richtungen der Koordinatenachsen $x,y,z$ durch die Punkte $(1)$, $(2)$, $(3)$, dann die Richtung der auf der Oberfläche senkrechten und in Bezug auf den Raumbereich $s$ nach außen gezogenen Normalen durch den Punkt $(4)$ und schließlich die Richtung der Geraden, die von einem beliebigen Punkt der Oberfläche zu seiner Lage nach der Variation gezogen wird, durch den Punkt $(5)$. Die Ortsänderung selbst, also die stets positiv zu nehmende Größe $\BareRoot{(\delta x^2+\delta y^2+\delta z^2)}$, bezeichnen wir der Kürze halber mit $\delta e$. Den Bogen zwischen zwei Punkten der Kugel, etwa $(1)$ und $(5)$, oder den Winkel, der diesen Bogen misst, schreiben wir als $(1,5)$. Damit gilt
\[
\delta x=\delta e.\cos(1,5),\quad \delta y=\delta e.\cos(2,5),\quad \delta z=\delta e.\cos(3,5)
\]

Dies gilt für jeden Punkt der Oberfläche. An ihrer Grenze, der Begrenzungslinie $P$, kommen zwei weitere Richtungen hinzu. Zuerst die Richtung des Elements $\mathrm{d}P$, der der
\SourceLeaf{68}Punkt $(6)$ entspreche, und dann die Richtung der zu diesem Element senkrechten Geraden, die die Oberfläche berührt und in Bezug auf sie nach innen gezogen wird. Ihr entspreche der Punkt $(7)$. Nach unserer Annahme liegen die Punkte $(6)$, $(7)$, $(4)$ in derselben Anordnung wie $(1)$, $(2)$, $(3)$. Zu beachten ist ferner, dass $(4,6)$, $(4,7)$, $(6,7)$ Quadranten oder rechte Winkel darstellen. Damit ergeben sich die bereits in Artikel 20 angegebenen Gleichungen
\[
\begin{gathered}
\eta Z-\zeta Y=\cos(1,7),\qquad \zeta X-\xi Z=\cos(2,7),\\
\xi Y-\eta X=\cos(3,7)
\end{gathered}
\]
und die Formeln des vorigen Artikels nehmen die folgenden Formen an.
\[
\begin{aligned}
Q&=-\delta e.\cos(5,7)\\
V&=\delta e.\cos(4,5).(\frac{\mathrm{d}\xi}{\mathrm{d}x}+\frac{\mathrm{d}\eta}{\mathrm{d}y})
\end{aligned}
\]
Somit gibt $Q$ für einen beliebigen Punkt der Begrenzungslinie $P$ die Verschiebung gegenüber der Ebene an, die diese Linie berührt und senkrecht zur Oberfläche $U$ steht. Sie wird auf der von dieser Oberfläche abgewandten Seite positiv gerechnet. Der Faktor von $V$, nämlich $\delta e.\cos(4,5)$, gibt offenbar die Verschiebung eines beliebigen Punkts der Oberfläche $U$ von ihrer Tangentialebene an, positiv gerechnet auf der vom Raumbereich $s$ abgewandten Seite.

Auch der andere Faktor von $V$ lässt sich durch seine geometrische Bedeutung erklären. Es gilt nämlich
\[
\begin{gathered}
\xi=-\zeta.\frac{\mathrm{d}z}{\mathrm{d}x},\qquad \eta=-\zeta.\frac{\mathrm{d}z}{\mathrm{d}y}\\
\frac{1}{\zeta\zeta}=1+(\frac{\mathrm{d}z}{\mathrm{d}x})^2+(\frac{\mathrm{d}z}{\mathrm{d}y})^2
\end{gathered}
\]
Daraus ergibt sich
\[
\begin{aligned}
\mathrm{d}\zeta&=\xi\zeta\zeta\mathrm{d}\frac{\mathrm{d}z}{\mathrm{d}x}+\eta\zeta\zeta\mathrm{d}\frac{\mathrm{d}z}{\mathrm{d}y}\\
\frac{\mathrm{d}\xi}{\mathrm{d}x}&=-\zeta\frac{\mathrm{dd}z}{\mathrm{d}x^2}-\frac{\mathrm{d}z}{\mathrm{d}x}.\frac{\mathrm{d}\zeta}{\mathrm{d}x}\\
&=-\zeta\frac{\mathrm{dd}z}{\mathrm{d}x^2}+\xi\xi\zeta\frac{\mathrm{dd}z}{\mathrm{d}x^2}+\xi\eta\zeta\frac{\mathrm{dd}z}{\mathrm{d}x.\mathrm{d}y}\\
&=-\zeta(\eta\eta+\zeta\zeta)\frac{\mathrm{dd}z}{\mathrm{d}x^2}+\xi\eta\zeta\frac{\mathrm{dd}z}{\mathrm{d}x.\mathrm{d}y}\\
\frac{\mathrm{d}\eta}{\mathrm{d}y}&=-\zeta\frac{\mathrm{dd}z}{\mathrm{d}y^2}+\eta\eta\zeta\frac{\mathrm{dd}z}{\mathrm{d}y^2}+\xi\eta\zeta\frac{\mathrm{dd}z}{\mathrm{d}x.\mathrm{d}y}\\
&=-\zeta(\xi\xi+\zeta\zeta)\frac{\mathrm{dd}z}{\mathrm{d}y^2}+\xi\eta\zeta\frac{\mathrm{dd}z}{\mathrm{d}x.\mathrm{d}y}
\end{aligned}
\]
und somit
\[
\begin{aligned}
\frac{\mathrm{d}\xi}{\mathrm{d}x}+\frac{\mathrm{d}\eta}{\mathrm{d}y}=-\zeta^3\big\{&\frac{\mathrm{dd}z}{\mathrm{d}x^2}[1+(\frac{\mathrm{d}z}{\mathrm{d}y})^2]-\frac{2\mathrm{dd}z}{\mathrm{d}x.\mathrm{d}y}.\frac{\mathrm{d}z}{\mathrm{d}x}.\frac{\mathrm{d}z}{\mathrm{d}y}\\
&+\frac{\mathrm{dd}z}{\mathrm{d}y^2}[1+(\frac{\mathrm{d}z}{\mathrm{d}x})^2]\big\}
\end{aligned}
\]
\SourceLeaf{69}Der Wert dieses Ausdrucks ist bekanntlich
\[
=\frac{1}{R}+\frac{1}{R'}
\]
wobei $R,R'$ die extremen Krümmungsradien im betrachteten Punkt bezeichnen und positiv zu rechnen sind, wenn die Oberfläche ihre konvexe Seite nach außen kehrt.

\Article{26}
Bei genauer Prüfung unserer Rechnung von Artikel 22 an wird eine stillschweigende Annahme sichtbar, die ihr zugrunde liegt. Beliebigen Werten der Koordinaten $x,y$ entspricht nur ein einziger Wert von $z$, und der Wert von $\zeta$ ist überall auf der ganzen Oberfläche $U$ positiv. Gleichwohl gilt der Satz, zu dem die Rechnung am Ende geführt hat, nämlich (I),
\[
\delta U=-\int\delta e.\cos(5,7).\mathrm{d}P+\int\delta e.\cos(4,5).(\frac{1}{R}+\frac{1}{R'})\mathrm{d}U
\]
allgemein und ist von dieser Annahme unabhängig. Hätten wir diese Allgemeinheit von Anfang an einbeziehen wollen, so hätten wir entweder einige Umständlichkeiten in Kauf nehmen oder ein etwas anderes Verfahren wählen müssen. Zum selben Ziel gelangen wir aber auch leicht durch die folgenden Überlegungen.

Unsere Rechnung ist offenbar unabhängig von der Annahme, dass die Koordinatenachse $z$ senkrecht steht. Vielmehr bleibt die Lage der Achsen darin völlig beliebig. Die Gültigkeit des Satzes steht für alle Oberflächen fest, für die sich die Gesamtheit der Punkte $(4)$ in eine einzige Halbkugel einschließen lässt. Es genügt, den Mittelpunkt (Pol) einer solchen Halbkugel als $(3)$ zu wählen.

Ist dagegen eine Oberfläche gegeben, die diese Bedingung nicht erfüllt, so lässt sie sich gewiss in zwei oder mehr Teile zerlegen, von denen jeder eine solche Bedingung erfüllt. Es ist leicht zu sehen, dass bei der Zerlegung einer Oberfläche in zwei Teile die Gültigkeit des Satzes für die ganze Figur unmittelbar aus seiner Gültigkeit für die einzelnen Teile folgt. Die Figur $U$ bestehe aus den Teilen $U',U''$, und $P'$ sei die Begrenzungslinie der Figur $U'$, $P''$ die der Figur $U''$. Ferner mögen $P',P''$ einen gemeinsamen Teil $P'''$ haben, sodass $P'$ aus $P'''$ und $P''''$ und $P''$ aus $P'''$ und $P'''''$ besteht. Offenbar hat dann die ganze Figur $U$ die Begrenzungslinie $P$, die sich aus $P''''$ und $P'''''$ zusammensetzt. Damit gilt zwar
\[
\begin{aligned}
\int\delta e.\cos(5,7)\mathrm{d}P'&=\int\delta e.\cos(5,7)\mathrm{d}P'''+\int\delta e.\cos(5,7)\mathrm{d}P''''\\
\int\delta e.\cos(5,7)\mathrm{d}P''&=\int\delta e.\cos(5,7)\mathrm{d}P'''+\int\delta e.\cos(5,7)\mathrm{d}P'''''
\end{aligned}
\]
\SourceLeaf{70}doch ist genau zu beachten, dass der Wert des Integrals $\int\delta e.\cos(5,7)\mathrm{d}P'''$, soweit es Teil des ersten Integrals ist, genau das Negative des Werts desselben Integrals ist, soweit es Teil des zweiten Integrals ist. Denn jedem Punkt der in beiden Fällen in entgegengesetzter Richtung durchlaufenen Linie $P'''$ entsprechen entgegengesetzte Lagen des Punkts $(7)$ und damit entgegengesetzte Werte des Faktors $\cos(5,7)$. Bei der Addition heben sich diese Teile also auf, und es ergibt sich
\[
\int\delta e.\cos(5,7)\mathrm{d}P'+\int\delta e.\cos(5,7)\mathrm{d}P''=\int\delta e.\cos(5,7)\mathrm{d}P
\]
Da $\delta U=\delta U'+\delta U''$ gilt, ergibt sich der Wert von $\delta U$ im Einklang mit der angegebenen Formel (I) unmittelbar, wenn angenommen wird, dass diese Formel mit den Werten der Variationen $\delta U',\delta U''$ übereinstimmt.

Schließlich bemerken wir, dass sich die Gültigkeit des Satzes (I) auch aus geometrischen Überlegungen hätte gewinnen lassen, und zwar leichter als auf dem analytischen Weg. Diesen haben wir hier dennoch vorgezogen, um die Gelegenheit zu nutzen, etwas Licht auf die Variationsrechnung für Doppelintegrale mit veränderlichen Grenzen zu werfen, die bisher wenig entwickelt ist. Das andere, geometrische Verfahren, das recht naheliegt, überlassen wir dem erfahrenen Leser.

\Article{27}
Es bleibt, die Variationen zu bestimmen, die die übrigen Elemente des Ausdrucks $W$ bei einer Änderung der Gestalt des Raumbereichs $s$ erfahren. Zuerst behandeln wir die Variation des Volumens des Raumbereichs $s$.

Nehmen wir die beiden in Artikel 21 betrachteten Dreiecke wieder auf und verbinden entsprechende Punkte ihrer Seiten durch Geraden, sodass ein Körper entsteht. An seine Stelle können wir ein Prisma mit der Grundfläche $\mathrm{d}U$ und der Höhe $\xi\delta x+\eta\delta y+\zeta\delta z=\delta e.\cos(4,5)$ setzen. Dieser Ausdruck liefert die Höhe positiv oder negativ, je nachdem, ob das verschobene Dreieck und damit der ganze Körper außerhalb oder innerhalb des Raumbereichs $s$ liegt. Damit haben wir (II)
\[
\delta s=\int\mathrm{d}U.\delta e.\cos(4,5)
\]
Weiter folgt daraus für die Variation des Integrals $\int z\mathrm{d}s$ die Formel (III)
\[
\delta\int z\mathrm{d}s=\int z\mathrm{d}U.\delta e.\cos(4,5)
\]

Was die Variation der Größe $T$ betrifft, so bemerken wir zunächst, dass $P$ die gemeinsame Grenze der Oberflächen $T,U$ bezeichnet. Die Verschiebungen der
\SourceLeaf{71}Punkte der Begrenzungslinie $P$ müssen also die Bedingung erfüllen, dass die neuen Lagen auf der Oberfläche des Raumbereichs $S$ bleiben. Durch die Verschiebung des Elements $\mathrm{d}P$ erfährt die Oberfläche $T$ also offenbar die Änderung $\pm\mathrm{d}P.\delta e.\sin(5,6)$. Man sieht leicht, dass das positive oder negative Vorzeichen im Allgemeinen vom Vorzeichen der Größe $\cos(4,5)$ abhängt. Diese Variation lässt sich aber übersichtlicher ausdrücken, wenn wir eine neue Richtung einführen. Sie liegt in der Tangentialebene an die Oberfläche des Raumbereichs $S$, steht senkrecht auf der Linie $P$ und weist in Bezug auf den Raumbereich $s$ nach außen. Bezeichnen wir den zu dieser Richtung gehörenden Punkt mit $(8)$, so beträgt die Variation der Oberfläche $T$, die durch die Verschiebung des Elements $\mathrm{d}P$ entsteht, $\mathrm{d}P.\delta e.\cos(5,8)$, oder (IV)
\[
\delta T=\int\mathrm{d}P.\delta e.\cos(5,8)
\]
wobei das Vorzeichen des Faktors $\cos(5,8)$ von selbst bestimmt, ob die Änderung eine Zunahme oder eine Abnahme ist.

Da der Punkt $(6)$ der Pol des durch die Punkte $(7)$, $(8)$ gezogenen Großkreises ist und der Punkt $(5)$ auf dem durch die Punkte $(6)$, $(8)$ gezogenen Großkreis liegt, bilden die Punkte $(5)$, $(7)$, $(8)$ ein bei $(8)$ rechtwinkliges Dreieck. Somit gilt $\cos(5,7)=\cos(5,8).\cos(7,8)$. Der Bogen $(7,8)$ misst aber den Winkel zwischen den beiden Ebenen, die die Oberflächen der Raumbereiche $s,S$ in ihrer Schnittlinie $P$ berühren, und zwar zwischen denjenigen Seiten dieser Ebenen, die den leeren Raum einschließen. Diesen Winkel bezeichnen wir mit $i$. Dann ist $180^0-i$ der Winkel zwischen den Seiten der Ebenen, die den Raumbereich $s$ enthalten, und unsere Formel (V) lautet
\[
\cos(5,7)=\cos(5,8).\cos i.
\]

